% Section added for thesis enhancement (Feb 2026)
\section{VS Code Extension: Interactive IDE Integration}
\label{sec:vscode_extension}

To bridge the gap between the backend LangGraph orchestrator and the developer's daily workflow, a custom Visual Studio Code extension was developed. This extension transforms the framework from a command-line tool into a real-time, context-aware assistant embedded directly in the Integrated Development Environment (IDE).

\subsection{Architecture Overview}

The integration follows a client-server architecture where the VS Code extension (TypeScript frontend) communicates with the FastAPI backend (Python) via HTTP streaming. This design provides several advantages:

\begin{itemize}
  \item \textbf{Technology Separation}: The complex LangGraph orchestration runs in Python, while the UI utilizes VS Code's native TypeScript APIs.
  \item \textbf{Asynchronous Execution}: Long-running workflows (NNI experiments, model downloads) don't freeze the VS Code interface.
  \item \textbf{Multi-Client Support}: The same backend can serve VS Code, CLI tools, and web interfaces simultaneously.
\end{itemize}

\subsection{Extension Configuration (\texttt{package.json})}

The extension manifest registers a Chat Participant using VS Code's native chat interface API (introduced in VS Code 1.85+).

\begin{lstlisting}[language=json, caption={Chat Participant registration in package.json}, label={lst:package_json}]
{
  "name": "stm32-ai-assistant",
  "displayName": "STM32 Edge AI Chat",
  "description": "Chat with your STM32 AI Workflow Agent",
  "version": "0.0.1",
  "engines": {
    "vscode": "^1.85.0"
  },
  "contributes": {
    "chatParticipants": [
      {
        "id": "stm32-ai.assistant",
        "name": "stm32ai",
        "fullName": "STM32 Edge AI Assistant",
        "description": "Analyze models, generate firmware, integrate AI",
        "isSticky": true
      }
    ]
  }
}
\end{lstlisting}

\textbf{Key Parameters:}
\begin{itemize}
  \item \texttt{id}: Unique identifier for the chat participant (used internally by VS Code).
  \item \texttt{name}: The \texttt{@stm32ai} handle users type to invoke the assistant.
  \item \texttt{isSticky: true}: Enables multi-turn conversations where context (previous messages) is automatically retained across chat interactions.
\end{itemize}

\subsection{TypeScript Handler Implementation}

The core logic resides in \texttt{extension.ts}, which implements a \texttt{ChatRequestHandler} callback.

\begin{lstlisting}[language=typescript, caption={Chat request handler in extension.ts}, label={lst:extension_handler}, mathescape=false]
const handler: vscode.ChatRequestHandler = async (
  request: vscode.ChatRequest,
  context: vscode.ChatContext,
  stream: vscode.ChatResponseStream,
  token: vscode.CancellationToken
) => {
  stream.markdown('Contacting STM32 Brain...\n\n');
  
  // 1. Prepare payload for FastAPI server
  const payload = {
    messages: [{ role: 'user', content: request.prompt }],
    user_id: "michele",
    session_id: "vscode-session"
  };
  
  // 2. POST request to streaming endpoint
  const response = await fetch('http://127.0.0.1:8000/stream', {
    method: 'POST',
    headers: { 'Content-Type': 'application/json' },
    body: JSON.stringify(payload)
  });
  
  // 3. NDJSON stream parsing
  const reader = response.body.getReader();
  const decoder = new TextDecoder("utf-8");
  let buffer = "";
  
  while (true) {
    if (token.isCancellationRequested) break;
    
    const { done, value } = await reader.read();
    if (done) break;
    
    buffer += decoder.decode(value, { stream: true });
    const lines = buffer.split("\n");
    buffer = lines.pop() || "";
    
    for (const line of lines) {
      if (!line.trim()) continue;
      const data = JSON.parse(line);
      
      if (data.type === 'markdown') {
        stream.markdown(data.content);
      } else if (data.type === 'progress') {
        stream.progress(data.content);
      } else if (data.type === 'error') {
        stream.markdown('**Error:** ' + data.content);
      }
    }
  }
};

vscode.chat.createChatParticipant('stm32-ai.assistant', handler);
\end{lstlisting}

\textbf{Implementation Details:}
\begin{itemize}
  \item \textbf{Cancellation Handling}: The \texttt{CancellationToken} allows users to cancel long-running requests by closing the chat window or clicking the stop button. The backend receives a dropped HTTP connection and gracefully terminates the LangGraph execution.
  \item \textbf{NDJSON Buffering}: Chunks may contain partial JSON objects. The \texttt{buffer} accumulates incomplete lines until a full newline-delimited JSON object is received.
  \item \textbf{Typed Responses}: The server emits different event types (\texttt{markdown}, \texttt{progress}, \texttt{error}), each rendered differently in the chat UI.
\end{itemize}

\subsection{FastAPI Backend: Server-Side Streaming}

The backend implements the \texttt{/stream} endpoint using FastAPI's \texttt{StreamingResponse}.

\begin{lstlisting}[language=python, caption={FastAPI streaming endpoint in server.py}, label={lst:server_streaming}]
@app.post("/stream")
async def stream_chat(request: ChatRequest):
    user_id = request.user_id
    session_id = request.session_id
    thread_id = f"{user_id}:{session_id}"
    
    # 1. Load user profile from Redis
    profile = await load_user_profile_from_redis(user_id)
    
    # 2. Build initial state
    initial_state = MasterState(
        message=request.messages[-1].content,
        last_board=profile.get("last_board"),
        preferred_compression=profile.get("preferred_compression")
    )
    
    # 3. Configure LangGraph thread
    config = {
        "configurable": {
            "thread_id": thread_id,
            "user_id": user_id
        }
    }
    
    # 4. Stream events from LangGraph execution
    async def event_generator():
        async for event in graph.astream(initial_state, config):
            if event["type"] == "node":
                yield f'{{"type": "progress", "content": "Executing {event["name"]}..."}}\n'
            elif event["type"] == "output":
                yield f'{{"type": "markdown", "content": {json.dumps(event["data"])}}}\n'
    
    return StreamingResponse(event_generator(), media_type="application/x-ndjson")
\end{lstlisting}

\textbf{Server-Side Features:}
\begin{itemize}
  \item \textbf{User Profile Injection}: The \texttt{last\_board} and \texttt{preferred\_compression} from Redis long-term memory pre-populate the state, enabling context chaining across sessions.
  \item \textbf{Thread-Based Isolation}: Each user+session combination gets a unique \texttt{thread\_id}, ensuring isolated checkpoints in Redis.
  \item \textbf{NDJSON Encoding}: Each event is serialized as a JSON object followed by \texttt{\textbackslash n}, enabling incremental parsing on the client.
\end{itemize}

\subsection{Communication Protocol}

The extension and server communicate via a custom NDJSON (Newline Delimited JSON) protocol. Each line represents a single event:

\begin{verbatim}
{"type": "progress", "content": "Routing to AI Flow..."}
{"type": "markdown", "content": "## Model Selected\n\nMobileNetV2 (0.35)"}
{"type": "progress", "content": "Downloading model from GitHub..."}
{"type": "markdown", "content": "Analysis complete. RAM: 280KB, Flash: 500KB"}
{"type": "error", "content": "Model exceeds STM32F4 Flash limit"}
\end{verbatim}

\textbf{Event Types:}
\begin{itemize}
  \item \texttt{progress}: Transient status updates (e.g., ``Analyzing model...''), rendered as bullet points in the chat.
  \item \texttt{markdown}: Rich-formatted responses supporting bold, code blocks, tables, and links.
  \item \texttt{error}: Critical failures displayed with red highlighting and error icons.
\end{itemize}

\subsection{Integration Challenges and Solutions}

\subsubsection{Challenge 1: LangGraph \texttt{config} Propagation}

\textbf{Problem:} During initial integration, the subgraph nodes did not receive the \texttt{config} dictionary, causing crashes when accessing \texttt{config["configurable"]["user\_id"]}.

\textbf{Solution:} Modified all node function signatures to include \texttt{config: dict = None} with default value \texttt{None}. Nodes check \texttt{if config and "configurable" in config} before accessing nested keys.

\subsubsection{Challenge 2: Infinite Clarification Loop}

\textbf{Problem:} Strict configuration validation caused the router to repeatedly request clarification (``Please specify target board'') even after the user provided input, creating a \texttt{router → clarify → router} loop.

\textbf{Solution:} Relaxed validation logic to allow \texttt{None} values for non-critical parameters. The router now proceeds with workflow execution and requests missing data via HITL interrupts later in the pipeline.

\subsubsection{Challenge 3: Model Download URL 404 Errors}

\textbf{Problem:} The STM32 ModelZoo GitHub repository URLs hardcoded in \texttt{predefined\_models.json} returned 404 errors due to repository restructuring.

\textbf{Solution:} Implemented fallback to GitHub search API when primary URL fails. The \texttt{search\_h5\_file\_in\_repo\_hybrid()} function scans repository contents recursively to locate valid \texttt{.h5} files.

\subsection{Developer Experience Benefits}

The VS Code integration provides several UX improvements over CLI-based interaction:

\begin{itemize}
  \item \textbf{Context Awareness}: The extension can access the active file (\texttt{vscode.window.activeTextEditor}), enabling future features like ``Analyze the model in this file'' or ``Generate firmware for this .ioc project.''
  \item \textbf{Inline Documentation}: Markdown responses support clickable file links (e.g., \texttt{file:///path/to/network.c\#L42}), allowing users to jump directly to generated code.
  \item \textbf{Progress Visibility}: Real-time progress updates reduce perceived latency for long-running operations (model downloads, NNI experiments).
  \item \textbf{Persistent Chat History}: Multi-turn conversations (\texttt{isSticky: true}) allow iterative refinement (``Actually, I want a smaller model'') without restarting the workflow.
\end{itemize}

\subsection{Deployment and Testing}

\textbf{Development Workflow:}
\begin{enumerate}
  \item \textbf{Compilation}: \texttt{npm run compile} transpiles TypeScript to JavaScript in \texttt{./out/extension.js}.
  \item \textbf{Debug Launch}: \texttt{F5} opens an Extension Development Host window where \texttt{@stm32ai} is available.
  \item \textbf{Server Startup}: The FastAPI backend must be running (\texttt{python -m src.api.server}) before testing.
\end{enumerate}

\textbf{Testing Scenarios:}
\begin{itemize}
  \item \textbf{Firmware Generation}: ``@stm32ai create a project for STM32H743 with STM32CubeIDE''
  \item \textbf{Model Analysis}: ``@stm32ai analyze MobileNetV2 for STM32F4 with high compression''
  \item \textbf{NNI Optimization}: ``@stm32ai fine-tune this model on Fruit-360 with 10 trials''
  \item \textbf{Cancellation Test}: Start a long-running task, then close the chat window to verify graceful termination.
\end{itemize}

\subsection{Future Enhancements}

Potential extensions to the VS Code integration:

\begin{itemize}
  \item \textbf{Syntax Highlighting for .ioc Files}: Custom TextMate grammar for STM32CubeMX configuration files.
  \item \textbf{Code Action Providers}: Right-click on an \texttt{.h5} model file → ``Analyze for STM32''
  \item \textbf{Status Bar Integration}: Show active NNI experiments in the VS Code status bar.
  \item \textbf{Terminal Integration}: Automatically open a terminal and flash firmware to connected STM32 boards via \texttt{st-flash}.
\end{itemize}
